\subsection{The actual global primitive comparison and exponential-gauge ratio}
\label{subsec:sr-source-exponential-gauge-ratio}

We retain here the literal algebraic global-functions law
used by the frozen source:
\[
 \Gamma(S,\mathcal O_S)\ \xrightarrow{\ p^*\ }\ 
 \Gamma(A,\mathcal O_A)\quad\hbox{is bijective
 for every abelian scheme }p:A\to S
 \hbox{ over every scheme }S.                               \tag{GR1}
\]
This is \(\mathrm{Out.GlobalFnsSt}\).
It is an algebraic all-base assertion; a theorem about
complex-valued functions on individual fibres does not
supply (GR1).  The conclusions below are its exact
applications in the chart comparison.

\begin{srlemma}[Normalized global generators are unique]
\label{lem:sr-source-normalized-global-generators-unique}
Let \(p:A\to S\) be an abelian scheme with identity \(o\)
and assume (GR1) for it.
Two global generators of the same actual line on \(A\)
which induce the same generator on \(o\) are equal.
This applies to the actual primitive generator and to
the following source line, for each \(a\in\mathcal O_E\):
\[
 {\cal L}_{\beta,a}
 =[1+a\bar a]^*L\otimes([a]^*L)^{-2}
                 \otimes L^{-1}\otimes[a\bar a]^*L^{-1}.
                                                               \tag{GR2}
\]
Its prescribed generator at the identity is
\(\mathrm{rig}^{-3}\).
\end{srlemma}
\begin{proof}
The ratio of the two line generators is a global unit
\(u\) on \(A\).
Construct it on each line chart by division of their
unit coefficients.  The transition coefficients cancel,
so the local unit ratios glue.
Their common normalization gives \(o^*u=1\), as an
identity of regular functions on \(S\).
By (GR1), \(u=p^*v\) for a regular function \(v\) on \(S\).
Since \(po=\mathrm{id}_S\), it follows that \(v=o^*u=1\),
and hence \(u=1\).
This uses no pointwise detection or reducedness assumption.
For (GR2), the exponent sum at the identity is
\(1-2-1-1=-3\); the normalization is therefore exactly
the displayed \(\mathrm{rig}^{-3}\).
The same ratio argument applies to the primitive line
\(L_xB(x,[w_0]x)^{-1}\) with its prescribed rigidification.
\end{proof}

Work now over any \(c.F\)-algebra \(S\), also a
commutative \(\mathbb Q\)-algebra.
Let \(x_m\) be one actual source exponential chart with
coordinates \(z_\lambda\) indexed by
\(\lambda\in\mathrm{Lt1}(D_w)\).
Write \(\tau_\lambda\) for the embedding underlying a
letter; its multiplicity index is retained separately.
Thus its point law is additive, its \(\mathcal O_E\)
action is diagonal,
\[
                    (D_a z)_\lambda=\chi_\lambda(a)z_\lambda,
 \qquad
 \chi_\lambda(a)=\mathrm{embK}_c(\tau_\lambda(a))\in S,             \tag{GR3}
\]
and its specified tangent points are the actual \(y.u_\lambda\).
An actual exponential gauge \(e=(e_m)\) is a family of
generators of \(x_m^*L\), compatible with truncation and
the origin rigidification, which also satisfies the
source's following global cocycle clause.
For every \(a\), there is a global generator
\(\beta_a\) of (GR2), normalized by
\(\mathrm{rig}^{-3}\), such that at every finite order
\[
 e_m(D_{1+a\bar a}z)
   =\beta_a(x_m(z))\,
            e_m(D_a z)^2 e_m(z)e_m(D_{a\bar a}z).            \tag{GR4}
\]
This notation denotes the actual tensor-line
\(\mathrm{Corr}\) identity; the four line factors and
their exponents are those of (GR2).
In particular, a different gauge may initially come
with a different global \(\beta_a\).

\begin{srproposition}[Classification of the ratio of two actual gauges]
\label{prop:sr-source-actual-exponential-gauge-ratio}
Assume (GR1).  Let \(e,e'\) be two actual exponential
gauges on this same actual exponential chart.
There is a formal series \(q\), independent of the
section subsequently evaluated, such that
\[
 \begin{gathered}
 q(0)=0,\qquad [z^\mu]q\ne0\ \Longrightarrow\ |\mu|=2,\\
 [z_\lambda z_{\lambda'}]q\ne0
             \ \Longrightarrow\ \tau_{\lambda'}=\overline{\tau_\lambda},\\
                         e'_m=(j_m\exp q)e_m.              \tag{GR5}
 \end{gathered}
\]
If an actual section of \(L^k\), \(k\ge0\), has
exponential germs \(g,g'\) in \(e,e'\), respectively,
then
\[
                              g'=g\exp(-kq).               \tag{GR6}
\]
Thus the literal \(\mathrm{GaugeUniqSt}\) follows from
the literal (GR1).
\end{srproposition}
\begin{proof}
Take the actual global unit ratio \(r_m\) of \(e'_m\)
to \(e_m\) on the affine jet scheme.
The two truncation identities imply
\(\rho_mr_{m+1}=r_m\): pull the ratio identity back,
compare both resulting \(\mathrm{Corr}\) identities,
and cancel the common unit generator on line charts.
Compatible finite polynomial representatives define a
formal series \(r\) with \(j_mr=r_m\), by the coefficient
construction used for (FC14).
The common rigidification gives \(r(0)=1\).

For each \(a\), the preceding lemma identifies the
two normalized global generators \(\beta_a,\beta'_a\).
Only after that global comparison may they be
cancelled from the two actual identities (GR4).
At each finite order this gives
\[
 r(D_{1+a\bar a}z)
                  =r(D_a z)^2r(z)r(D_{a\bar a}z).           \tag{GR7}
\]
It holds as a formal identity since it holds at every
jet order.
This cancellation is the source's use of global
primitive-line data; local scalar comparison alone
would not give (GR7).

Let \(q=\log r\), using the formal logarithm of a
unit with constant term \(1\).
Formal exponential and logarithm are inverse on
zero-constant and constant-one series over any
\(\mathbb Q\)-algebra.
Their identities can be verified in each nilpotent
jet quotient, where the sums are finite; no analytic
convergence is involved.
Taking logarithms of (GR7) gives
\[
 q(D_{1+a\bar a}z)
              =2q(D_a z)+q(z)+q(D_{a\bar a}z).              \tag{GR8}
\]
Take \(a=1\).  At a degree-\(j\) coefficient \(c\),
this becomes \((2^j-4)c=0\).
For \(j\ne2\), \(2^j-4\) is a nonzero rational and
therefore a unit in \(S\).
Thus \(q\) is homogeneous of degree two.

We include the source's character calculation, since
cancellation of a nonzero coefficient is not available
over the allowed rings.
For a coefficient \(c=[z_\lambda z_{\lambda'}]q\), set
\(s=\chi_\lambda\), \(t=\chi_{\lambda'}\).
Equation (GR8) at that coefficient is exactly
\[
 c\{s(a)s(\bar a)+t(a)t(\bar a)-2s(a)t(a)\}=0
                               \qquad(a\in\mathcal O_E).   \tag{GR9}
\]
Write \(F(a)\) for its left side.
Use (GR9) at \(a,a+1,a^2,a^2+1\).
Expanding these four identities gives
\[
 \begin{split}
 2c(s(a)-s(\bar a))(s(a)-t(\bar a))
  ={}&-\{F(a^2+1)-F(a^2)\}\\
    &+\{t(a)+s(\bar a)+t(\bar a)-s(a)\}
                                      \{F(a+1)-F(a)\}.
 \end{split}                                                   \tag{GR10}
\]
The right side is zero; cancel \(2\), a unit.

Here is the exact unit-separation inference.
Suppose characters \(\sigma,\tau_1,\tau_2\) satisfy
\[
        c(\sigma(a)-\tau_1(a))(\sigma(a)-\tau_2(a))=0
                                                    \quad\forall a
                                                               \tag{GR11}
\]
and distinct character indices can be separated by a
unit difference at an integral element.
If \(\sigma\ne\tau_1,\tau_2\), choose \(a_1\) separating
the first pair and \(a_2\) separating the second.
At \(a_1\) cancel the first unit to get
\(c(\sigma(a_1)-\tau_2(a_1))=0\).
At \(a_2\) cancel the second to get
\(c(\sigma(a_2)-\tau_1(a_2))=0\).
Expand (GR11) at \(a_1+a_2\); these two identities
remove every term except
\[
             c(\sigma(a_1)-\tau_1(a_1))
                       (\sigma(a_2)-\tau_2(a_2))=0.
\]
Both factors are units, so \(c=0\).
Consequently \(c\ne0\) forces
\(\sigma=\tau_1\) or \(\sigma=\tau_2\).

Apply this to (GR10) and also with \(s,t\) interchanged.
It gives, for \(c\ne0\),
\[
 s=\bar s\ \hbox{or}\ s=\bar t,\qquad
 t=\bar t\ \hbox{or}\ t=\bar s.                             \tag{GR12}
\]
If either second alternative holds, conjugate
involutivity gives \(t=\bar s\).
In the remaining case both characters are fixed by
conjugation.  Equation (GR9) then becomes
\(c(s(a)-t(a))^2=0\).
A unit difference forces \(c=0\) unless \(s=t\);
when \(s=t\), again \(t=\bar s\).
This also covers the general character version of the
source lemma, before specializing to CM embeddings.

For our actual embeddings, the required unit separation
is precisely the integral embedding separation in
(FC16): a nonzero difference in the field \(c.F\)
maps to a unit in \(S\).
For embedding characters,
\(\chi_{\bar\tau}(a)=\chi_\tau(\bar a)\),
and conjugation is involutive.
The conclusion compares the underlying embeddings
\(\tau_\lambda,\tau_{\lambda'}\), without identifying
their multiplicity indices.
The character conclusion therefore proves the
support assertion in (GR5).
Since \(r=\exp q\), its frame-ratio assertion follows.

Finally take actual section-pullback and value
witnesses for \(g,g'\) at order \(m\).
The two pulled sections are equal by uniqueness.
Their generators differ by \(r_m^k\), so cancelling
the original unit generator gives
\[
                           j_mg=r_m^k j_mg'.
\]
Multiply by \(j_m\exp(-kq)\), using
\((j_m\exp q)^k j_m\exp(-kq)=1\).
This proves (GR6) at every order.
Coefficientwise equality of all jets proves the
formal-series identity (GR6).
The same \(q\) works for all \(k\) and all actual
sections because it was constructed from \(e,e'\)
before they were chosen.
\end{proof}

\paragraph{Frozen source correspondence and supplier boundary.}
The exact all-base law (GR1) is
\(\mathrm{Out.GlobalFnsSt}\) in
\texttt{src/ReciprocityTheory/ReductiveGroups/ArithmeticGeometry/}
\texttt{ExtensionDescent/GlobalSectionsRankAlgClosed.lean}
(\texttt{372ffae0095db579e74998eb5a086acc55e14f8d}).
The actual definitions (GR2)--(GR4), including
\(\mathrm{beta\varepsilon}=(1,-2,-1,-1)\) and the
\(-3\) rigidification, are
\texttt{src/Formalism/LocalToGlobal/ProjectiveLattices/IsExpChart.lean}
(\texttt{f352d7a79e14580fbb6b765c72e4fddf814f8fee}).
All following files have prefix
\texttt{src/Formalism/SharpRestriction/Divisibility/GeneratorPowers/}.
The global ratio and primitive uniqueness are
\texttt{PrimUniqPf.lean}
(\texttt{4b27f88033c458bcffadb51a77f861dd5d50d2d2});
the exact \(\beta\)-uniqueness is
\texttt{GaugeBeta.lean}
(\texttt{a014cc8f5289738ff9985fb0e463d74e8f9a6b95}).

The unit-ratio comparisons, inverse-limit coefficients,
embedding-character transport, and actual section-value
scaling are respectively
\texttt{GaugeUniqA.lean}
(\texttt{525d54081edda65391bbe14f004140f886f0ecfa}),
\texttt{GaugeUniqB.lean}
(\texttt{83424abbfb926b029dccf4bd0de7acde1d384cc3}),
\texttt{GaugeUniqC.lean}
(\texttt{2b513db9858e2f312fafb6b8e2341568dd196b6a}),
and \texttt{GaugeUniqD.lean}
(\texttt{b0957cbea0258f7f987166ce3b6be6b2c912959c}).
The full actual assembly is
\texttt{GaugeUniqPf.lean}
(\texttt{9eec11ab88ee39dba55f932b8300f80bacbdfc0c}),
\(\mathrm{gaugeUniq\_of\_global}\).
The logarithm, degree-two and character calculation
(GR8)--(GR12) are the actual \(\mathrm{gaugeRatio}\)
proof in
\texttt{src/ReciprocityTheory/ReductiveGroups/ArithmeticGeometry/}
\texttt{TaylorDiscriminants/WickVanishing/GaugeRatio.lean}
(\texttt{e2bed2ce5f561059bab0979ac1d3c276803a8e90}).
In particular, the unit-separated calculation is retained;
a domain-only character argument would give a weaker result.

The present section constructs the gauge-comparison
conclusion from (GR1) for two supplied actual
exponential gauges.  It does not construct the gauges
from the original chart and its canonical formal
connection, or identify two actual exponential chart
point systems.  Those producers and harmonic
annihilation are required before this result supplies
arbitrary-chart independence.
